% !TEX root = ../../main.tex
%==================================================================
% chapters/theory/mc5_genus0_genus1_wall_platonic.tex
%
% MC5: Five-point sewing and genus-one clutching.
%
% Purpose: inscribe the 5-point sewing theorem of MC5.  The rational
% compactification \overline M_{0,5} contributes its ten boundary
% divisors; the first genus-one datum enters through the clutching
% correspondence in \overline M_{1,5}.
%
% Cross-references: chap:climax-platonic (eq:climax-equation-G2 supplies
% the bar differential as KZ pullback), chap:universal-conductor (kappa
% conductor identity), chap:shadow-quadrichotomy-platonic (G/L/C/M
% partition controlling which sewing prescription instantiates),
% chap:mc5-class-m-chain-level-platonic (chain-level pro-ambient
% companion theorem).
%==================================================================

\chapter{MC5: Five-point sewing and genus-one clutching}
\label{chap:mc5-g0g1-wall-platonic}

% Local macros (idempotent providecommand; required because chapter
% files do not introduce new macros, AP-LATEX). Names are namespaced
% with mc5g0g1 to avoid collision with the chain-level companion.
\providecommand{\KZ}{\mathrm{KZ}}
\providecommand{\Arn}{\mathrm{Arnold}}
\providecommand{\dbar}{d_{\mathrm{bar}}}
\providecommand{\MMod}{\overline{M}}
\providecommand{\Mbar}{\overline{M}}
\providecommand{\Conf}{\mathrm{Conf}}
\providecommand{\Conford}{\mathrm{Conf}^{\mathrm{ord}}}
\providecommand{\BarOrd}{\bar B^{\mathrm{ch}}}
\providecommand{\Conv}{\mathrm{Conv}}
\providecommand{\MC}{\mathrm{MC}}
\providecommand{\AInf}{A_{\infty}}
\providecommand{\evcore}{\mathrm{ec}}
\providecommand{\fg}{\mathfrak{g}}

The master conjecture chain $\mathrm{MC1} \to \mathrm{MC2} \to \mathrm{MC3} \to
\mathrm{MC4}$ closes inside genus zero. The four-point function lives on
$\Mbar_{0,4} \cong \mathbb{P}^1$; its three boundary points are rational
nodal degenerations corresponding to the three pairings $\{12|34\},
\{13|24\}, \{14|23\}$ of the four marked points. No genus-1 stratum
appears, and the entire MC4 sewing closure proceeds inside genus zero.

The five-point function lives on $\Mbar_{0,5}$. Its Deligne--Mumford
boundary has only rational separating divisors: ten of them, indexed by
two-element subsets of the five markings. The first genus-one datum is a
different geometric operation. It is the clutching correspondence
\[
\mathrm{cl}_{T}\colon
\Mbar_{1,T\cup\{*\}}\times \Mbar_{0,T^{c}\cup\{*\}}
\longrightarrow \Mbar_{1,5},
\qquad |T|=3,
\]
where the node labelled $*$ glues a genus-one component carrying the
three markings in~$T$ to a residual genus-zero component carrying the
remaining two. This correspondence, not an extra divisor of
$\Mbar_{0,5}$, is the genus-zero/genus-one wall seen by MC5.

This chapter inscribes the 5-point sewing theorem in that form. The
statement has five clauses: rational separating divisors, genus-one
clutching charts, pullback identity from the climax theorem, completed
Maurer--Cartan element, and bar-cobar duality. The
evaluation-generated-core qualifier is preserved as a mathematical
hypothesis and complemented by a full-category conjecture. The
chain-level companion
\textup{(}Chapter~\ref{chap:mc5-class-m-chain-level-platonic}\textup{)}
addresses the ambient-completion question for class $\mathsf{M}$; the
present chapter addresses the geometric sewing content.

\section{Setup: the moduli geometry of \texorpdfstring{$\Mbar_{0,5}$}{Mbar(0,5)}}
\label{sec:mc5-g0g1-setup}

Throughout this chapter, $X$ is a smooth projective curve over
$\mathbb{C}$, and $\cA$ is a strong completion tower
\textup{(}Definition~\ref{def:strong-completion-tower}\textup{)} of finite
resonance rank $\rho(\cA) < \infty$. The evaluation-generated subcategory
$\cA^{\evcore}$ is defined as in MC3
\textup{(}Section~\ref{sec:mc3-platonic-scope}\textup{)}.

\begin{definition}[Rational boundary and genus-one clutching]
\label{def:mbar-05-boundary-g0g1-wall}
The Deligne--Mumford compactification $\Mbar_{0,5}$ of the moduli of
$5$-pointed genus-zero stable curves has a normal-crossings boundary
divisor that decomposes as
\begin{equation}
\label{eq:mbar-05-boundary-decomposition}
\partial \Mbar_{0,5}
\;=\;
\bigcup_{\substack{S \subset \{1,\dots,5\} \\ |S| = 2}}
D_{S},
\qquad
D_{S}=\delta_{0|S,S^c}.
\end{equation}
Here $\delta_{0|S, S^c}$ is the locus of stable curves with two
genus-zero components attached at a single node, the markings indexed by
$S$ on one component and the markings indexed by $S^c$ on the other.
There are $\binom{5}{2}=10$ such divisors.

For a three-element subset $T\subset\{1,\ldots,5\}$, the genus-one
sewing datum is the clutching morphism
\begin{equation}\label{eq:mc5-genus-one-clutching}
\mathrm{cl}_{T}\colon
\Mbar_{1,T\cup\{*\}}\times \Mbar_{0,T^{c}\cup\{*\}}
\longrightarrow \Mbar_{1,5}.
\end{equation}
Its image $\Delta^{(1)}_{T}$ is a boundary divisor of
$\Mbar_{1,5}$, not of $\Mbar_{0,5}$. The elliptic trace in the MC5
sewing theorem is a compatibility condition along the completed charts
of~$\Delta^{(1)}_{T}$.
\end{definition}

The decomposition~\eqref{eq:mbar-05-boundary-decomposition} contains ten
rational divisors. The genus-one charts
\eqref{eq:mc5-genus-one-clutching} are separate clutching
correspondences. The point counts visible to the MC chain at each prior
stage are summarised in the following table.

\begin{table}[t]
\centering
\begin{tabular}{l@{\quad}l@{\quad}l}
\toprule
Stage & Moduli accessed & Genus filter \\
\midrule
$\mathrm{MC1}$ & $\Mbar_{0,1}$ & $g = 0$ \\
$\mathrm{MC2}$ & $\Mbar_{0,2}$ & $g = 0$ \\
$\mathrm{MC3}$ & $\Mbar_{0,3}$ & $g = 0$ \\
$\mathrm{MC4}$ & $\Mbar_{0,4} \cong \mathbb{P}^{1}$ & $g = 0$ \\
$\mathrm{MC5}$ & $\Mbar_{0,5}$ plus $\mathrm{cl}_{T}$ for $|T|=3$ & $g = 0$ with first $g=1$ clutching \\
\bottomrule
\end{tabular}
\caption{The MC chain as a stratified arithmetic of point counts. MC5
is the first stage at which the genus-one clutching datum is visible.}
\label{tab:mc-chain-genus-filter-g0g1-wall}
\end{table}

\begin{remark}[The arithmetic of MC1--MC5]
\label{rem:mc5-arithmetic-g0g1-wall}
At MC1 the moduli is a point, the MC element is the unit-vacuum
identification, and no boundary structure appears. At MC2 the moduli is
the line $\Mbar_{0,2}$, the MC element is the binary OPE residue at the
collision divisor (the chiral classical $r$-matrix), and the only
boundary geometry is the single nodal point. At MC3 the moduli is
$\Mbar_{0,3}$, a single point, and the MC element encodes the first
associativity homotopy at the level of the bar coderivation $b_{2}+b_{3}$
\textup{(}see Theorem~\ref{thm:mc3-evaluation-core-five-family}\textup{)}. At MC4 the
moduli $\Mbar_{0,4} \cong \mathbb{P}^{1}$ has three boundary points, all
rational, and the strong-completion-tower closure of
Theorem~\ref{thm:completed-bar-cobar-strong} extends the MC element
across them. At MC5 the moduli $\Mbar_{0,5}$ has ten rational
separating divisors, and the genus-one trace enters through the
clutching maps~\eqref{eq:mc5-genus-one-clutching}.
\end{remark}

\section{The 5-point sewing theorem}
\label{sec:mc5-statement-g0g1-wall}

The central theorem of the chapter, stated now, is restricted to the
evaluation-generated core $\cA^{\evcore}$ in keeping with the
eval-core discipline (thick generation on evaluation-generated
representations, not full $\DK_g$) applied throughout MC3 and MC4
\textup{(}see Remark~\ref{rem:mc5-g0g1-evcore-qualifier} below for the
explicit scope statement\textup{)}.

\begin{theorem}[MC5 5-point sewing with genus-one clutching;
\ClaimStatusProvedHere]
\label{thm:mc5-g0g1-wall-five-point-sewing}
\index{MC5!5-point sewing|textbf}
\index{genus-one clutching|textbf}
Let $X$ be a smooth projective curve over $\mathbb{C}$ and let $\cA$ be a
strong completion tower of finite resonance rank $\rho(\cA) < \infty$,
restricted to its evaluation-generated core $\cA^{\evcore}$.
For $\xi_{1}, \dots, \xi_{5} \in \cA^{\evcore}$ and $z_{1}, \dots, z_{5}
\in X$ with $z_{i} \neq z_{j}$ for $i \neq j$, let
\[
\Phi^{\evcore}_{5}(\xi_{1}, \dots, \xi_{5}; z_{1}, \dots, z_{5})
\;=\;
\bigl\langle \mathrm{vac} \,\bigm|\, V(\xi_{1}, z_{1}) \cdots V(\xi_{5}, z_{5}) \,\bigm|\, \mathrm{vac}\bigr\rangle
\]
be the $5$-point chiral correlator of $\cA^{\evcore}$ on $X$, regarded
as a holomorphic section of the conformal-weight line over
$\Conford_{5}(X)$. Then $\Phi^{\evcore}_{5}$ has a unique rational
sewing extension over the completed genus-zero compactification
$\widehat{\Mbar}_{0,5}(X)$. For each three-element subset
$T\subset\{1,\ldots,5\}$ there is also a genus-one sewn block
$\bar\Phi^{(1),\evcore}_{5,T}$ on the completed clutching chart
$\widehat{\Delta}^{(1)}_{T}$, compatible with the rational nodal
factorisation at $q=0$. The sewing data are:
\begin{enumerate}[label=\textup{(\roman*)}]
\item \textbf{(Genus-0 separating divisors.)} For each separating divisor
$\delta_{0|S, S^c}$ with $|S| \in \{2, 3\}$, the restriction
\begin{equation}
\label{eq:mc5-g0g1-restriction-separating}
\bar\Phi^{\evcore}_{5}\big|_{\delta_{0|S, S^c}}
\;=\;
\sum_{a, b}
\Phi^{\evcore}_{|S|+1}\bigl(\xi_{S}; e_{a}\bigr) \,
\eta^{ab} \,
\Phi^{\evcore}_{|S^c|+1}\bigl(e_{b}; \xi_{S^c}\bigr)
\end{equation}
agrees with the MC4 \textup{(}respectively MC3\textup{)} sewing
prediction, where $\{e_{a}\}$ is a homogeneous basis of the
eval-gen-core sectors with parameters $a, b$ ranging over a fixed
indexing set, and $\eta^{ab}$ is the inverse Shapovalov form on the node
sector. The sum is convergent and trace-class on the appropriate
weighted Hilbert completion.

\item \textbf{(Genus-one clutching.)} For each
$T=\{i_{1},i_{2},i_{3}\}$ with complement
$T^c=\{j_{1},j_{2}\}$, along the completed clutching chart
$\widehat{\Delta}^{(1)}_{T}$ parametrised by the bubble nome
$q \in \{q \in \mathbb{C}:0<|q|<1\}$,
the restriction
\begin{equation}
\label{eq:mc5-g0g1-restriction-elliptic-corner}
\bar\Phi^{(1),\evcore}_{5,T}
\;=\;
\sum_{a, b}
\mathrm{Tr}^{(1)}_{\cA^{\evcore}}\Bigl(V(\xi_{i_{1}}) V(\xi_{i_{2}}) V(\xi_{i_{3}}) \, e_{a} \, q^{L_{0} - c/24}\Bigr)
\,\eta^{ab}\,
\Phi^{\evcore}_{2}(\xi_{j_{1}}, \xi_{j_{2}}; e_{b})
\end{equation}
is the elliptic supertrace of the $3$-point function on the genus-one
component glued through the node state $e_a$ to the residual
$2$-point function on the genus-zero component.

\item \textbf{(Pullback identity from the climax theorem.)} The bar
differential $d_{\mathrm{bar}, 5}$ of $\BarOrd_{\le 5}(\cA^{\evcore})$
is the pullback
\begin{equation}
\label{eq:mc5-g0g1-bar-differential-as-pullback}
d_{\mathrm{bar}, 5} \;=\; \KZ^{*}\bigl(\nabla^{(5)}_{\Arn}\bigr)
\end{equation}
of the $n = 5$ specialisation of Arnold's universal flat connection on
$\Conford_{5}(X)$ along the KZ functor of
Chapter~\ref{chap:climax-platonic}, equation~\eqref{eq:climax-equation-G2}.
The Arnold flatness identity $\nabla^{(5)\, 2}_{\Arn} = 0$ is, after
pullback, the bar squaring-to-zero $d_{\mathrm{bar}, 5}^{\, 2} = 0$.

\item \textbf{(Completed Maurer-Cartan element.)} The sewing twisting
morphism
\begin{equation}
\label{eq:mc5-g0g1-completed-mc-element}
\tau^{\mathrm{sew}, 5}
\;\in\;
\MC\Bigl(\Conv\bigl(\BarOrd_{\le 5}(\cA^{\evcore}), \cA^{\evcore}\bigr)\Bigr)
\end{equation}
extends the MC4 element $\tau^{\mathrm{sew}, 4}$, agrees with the
genus-$0$ separating prediction~\eqref{eq:mc5-g0g1-restriction-separating}
along each $\delta_{0|S, S^c}$ for $|S| \in \{2, 3\}$, and agrees with
the genus-one clutching prediction
\eqref{eq:mc5-g0g1-restriction-elliptic-corner} after applying the
clutching functor for every $T$ with $|T|=3$.

\item \textbf{(Bar-cobar duality.)} The completed counit
\begin{equation}
\label{eq:mc5-g0g1-bar-cobar-counit}
\widehat{\epsilon}_{5} \;\colon\;
\widehat{\Omega}^{\mathrm{ch}}\bigl(\BarOrd_{\le 5}(\cA^{\evcore})\bigr)
\xrightarrow{\;\sim\;}
\cA^{\evcore}_{\le 5}
\end{equation}
is a quasi-isomorphism of curved dg chiral algebras at finite stage
$N = 5$, and its degree-$5$ component coincides with the cobar dual of
$\tau^{\mathrm{sew}, 5}$.
\end{enumerate}
\end{theorem}

\begin{remark}[evaluation-generated-core qualifier]
\label{rem:mc5-g0g1-evcore-qualifier}
Theorem~\ref{thm:mc5-g0g1-wall-five-point-sewing} is proved on the
evaluation-generated subcategory $\cA^{\evcore}$, exactly in the sense of
MC3 \textup{(}see Remark~\ref{rem:mc3-eval-generated-core}\textup{)}.
The full-category extension to objects outside the eval-gen-core is
conjectural for two structural reasons. First, the elliptic supertrace
in clause~\textup{(ii)} requires trace-class compactness of the
restriction operator on the genus-$1$ bubble, which is proved by the
HS-sewing criterion of Theorem~\ref{thm:hs-sewing-trace-class} only on
eval-gen-core objects: the OPE coefficient growth bound is verified
family-by-family for the eval-gen-core, not for arbitrary chiral algebras.
Second, the bar-cobar quasi-isomorphism in clause~\textup{(v)} inherits
its eval-gen-core restriction from the corresponding MC4 statement
\textup{(}Theorem~\ref{thm:completed-bar-cobar-strong}\textup{)} via the
finite-type bar-cobar comparison. Stating
Theorem~\ref{thm:mc5-g0g1-wall-five-point-sewing} without the eval-gen-core
qualifier would extend the claim to the full category $\DK_g$, where
the OPE-coefficient-growth bound and the MC4 bar-cobar comparison have
not been established; the eval-core discipline (thick generation on
evaluation-generated representations) is therefore mandatory throughout
the MC chain.
\end{remark}

\begin{conjecture}[Full-category MC5 with genus-one clutching]
\label{conj:mc5-g0g1-wall-full-category}
\index{MC5!full-category, conjectural at g0/g1 wall|textbf}
The conclusions of Theorem~\ref{thm:mc5-g0g1-wall-five-point-sewing}
extend to the full subcategory of pronilpotent strong completion towers
of finite resonance rank, without the evaluation-generated-core
restriction.
\ClaimStatusConjectured{}
\end{conjecture}

\begin{remark}[Higher-stage MC extension is conjectural]
\label{rem:mc5-g0g1-higher-stage-conjectural}
The MC chain $\mathrm{MC1} \to \cdots \to \mathrm{MC5}$ is closed at
$\mathrm{MC5}$ on the eval-gen-core. Extension to $\mathrm{MC6}$ and
beyond would force additional genus-filter crossings: at $n = 6$ markings,
either two disjoint genus-one clutching components appear or the first
genus-two clutching component is reached. The corresponding sewing
theorems are conjectural; their proofs would require trace-class
compactness on the higher-genus component, beyond the HS-sewing
criterion currently available on the eval-gen-core
\textup{(}with~\ref{conj:mc5-g0g1-wall-full-category} the simplest
analogue\textup{)}.
\end{remark}

\section{Stratified extension and genus-one clutching}
\label{sec:mc5-g0g1-stratified-extension}

The proof of Theorem~\ref{thm:mc5-g0g1-wall-five-point-sewing} proceeds
by extension across the ten rational boundary divisors of $\Mbar_{0,5}$
and by a separate compatibility check on the genus-one clutching
charts~\eqref{eq:mc5-genus-one-clutching}. The rational divisors are
governed by MC4 and MC3; the clutching chart is governed by the
genus-one trace formula on the elliptic component.

\begin{proof}[Proof of Theorem~\ref{thm:mc5-g0g1-wall-five-point-sewing}]
The proof has four steps corresponding to the four sewing clauses.
Throughout, $\cA$ is restricted to $\cA^{\evcore}$
\textup{(}see Remark~\ref{rem:mc5-g0g1-evcore-qualifier}\textup{)}.

\medskip
\textbf{Step 1: Restriction to rational boundary divisors and clutching charts.}
The boundary $\partial \Mbar_{0,5}$ is decomposed as
in~\eqref{eq:mbar-05-boundary-decomposition}: ten rational separating
divisors $D_S$ with $|S|=2$. The extension
$\bar\Phi^{\evcore}_{5}$ is determined by its restriction to these
divisors via the standard formal-coordinate factorisation: write
$z_i=z_i(t)$ with $t$ the formal parameter normal to the divisor, expand
$\Phi^{\evcore}_{5}$ as a power series in $t$, and apply the OPE
expansion encoded by clause~(i). Clause~(ii) is not another
$\Mbar_{0,5}$ boundary condition; it is checked on the completed
clutching chart $\widehat{\Delta}^{(1)}_{T}\subset\Mbar_{1,5}$.

\medskip
\textbf{Step 2: Compatibility with MC4 at separating divisors.}
At each separating divisor $\delta_{0|S, S^c}$ with $|S| = 2$, the
curve degenerates into two genus-$0$ components glued at a single node;
the two markings $\xi_{S}$ live on one component (the bubble), and the
three markings $\xi_{S^c}$ live on the other. The factorisation through
the node is given by the complete intermediate state sum
\begin{equation}
\label{eq:mc5-g0g1-mc4-intermediate-sum}
\Phi^{\evcore}_{5}\big|_{\delta_{0|S, S^c}}
\;=\;
\sum_{a, b}
\Phi^{\evcore}_{3}\bigl(\xi_{S}; e_{a}\bigr)
\,\eta^{ab}\,
\Phi^{\evcore}_{4}\bigl(e_{b}; \xi_{S^c}\bigr),
\end{equation}
where the sum runs over the eval-gen-core basis $\{e_{a}\}$. The
right-hand side is exactly the MC3-MC4 sewing prediction: the
$3$-point factor is governed by MC3
\textup{(}Theorem~\ref{thm:mc3-evaluation-core-five-family}\textup{)}, and the $4$-point
factor is governed by MC4
\textup{(}Theorem~\ref{thm:completed-bar-cobar-strong}\textup{)},
both restricted to $\cA^{\evcore}$ as in
Remark~\ref{rem:mc5-g0g1-evcore-qualifier}.

The case $|S| = 3$ is symmetric: the bubble carries three markings, and
the residual sphere carries two. Expanded out,
\begin{equation}
\label{eq:mc5-g0g1-mc4-intermediate-sum-symmetric}
\Phi^{\evcore}_{5}\big|_{\delta_{0|S, S^c}}
\;=\;
\sum_{a, b}
\Phi^{\evcore}_{4}\bigl(\xi_{S}; e_{a}\bigr)
\,\eta^{ab}\,
\Phi^{\evcore}_{3}\bigl(e_{b}; \xi_{S^c}\bigr),
\end{equation}
with the symmetric MC4 closure on the bubble factor and the MC3 closure
on the residual factor.

Convergence and trace-class compactness of the sums in
\eqref{eq:mc5-g0g1-mc4-intermediate-sum} and
\eqref{eq:mc5-g0g1-mc4-intermediate-sum-symmetric} follow from the
HS-sewing criterion of Theorem~\ref{thm:hs-sewing-trace-class}, applied
at $n = 4$ and at $n = 3$ respectively. The strong-completion-tower
hypothesis of $\cA^{\evcore}$ controls the OPE coefficient growth
required by the criterion.

\medskip
\textbf{Step 3: Compatibility on the genus-one clutching chart.}
Fix $T=\{i_1,i_2,i_3\}$. The chart
$\widehat{\Delta}^{(1)}_{T}$ is the formal neighbourhood of the image of
\eqref{eq:mc5-genus-one-clutching} in $\Mbar_{1,5}$. Its smoothing
parameter is the bubble nome $q=e^{2\pi i\tau}$ with $0<|q|<1$; the
limit $q\to0$ pinches the elliptic component to its nodal
degeneration.

On the bubble, the $3$-point function is determined by the genus-$1$
traced $3$-point function
\begin{equation}
\label{eq:mc5-g0g1-genus-1-traced-three-point}
F^{(1)}_{3}(\xi_{i_{1}}, \xi_{i_{2}}, \xi_{i_{3}}; e_{a}; q)
\;=\;
\mathrm{Tr}_{\cA^{\evcore}}\Bigl(V(\xi_{i_{1}}) V(\xi_{i_{2}}) V(\xi_{i_{3}})
\, e_{a} \, q^{L_{0} - c/24}\Bigr),
\end{equation}
where the trace is over the eval-gen-core sectors. By the
factorisation-at-nodes clause of Theorem~A
\textup{(}Theorem~\ref{thm:theorem-A-modular-family}, clause~(iv)\textup{)},
the bar differential on a degenerating curve restricts at the node to
the ordinary residue OPE, so
\begin{equation}
\label{eq:mc5-g0g1-bar-differential-restriction-corner}
d_{\mathrm{bar},T}^{(1|0)}
\;=\;
d^{(1)}_{\mathrm{bar}, 3} \otimes 1 \;+\; 1 \otimes d^{(0)}_{\mathrm{bar}, 2},
\end{equation}
the formal sum of the genus-$1$ bar differential on the bubble with the
genus-$0$ bar differential on the residual sphere. The $5$-point sewing
prediction on the clutching chart is therefore the trace formula
\begin{equation}
\label{eq:mc5-g0g1-corner-trace-formula}
\bar\Phi^{(1),\evcore}_{5,T}
\;=\;
\sum_{a, b}
F^{(1)}_{3}(\xi_{i_{1}}, \xi_{i_{2}}, \xi_{i_{3}}; e_{a}; q)
\,\eta^{ab}\,
\Phi^{\evcore}_{2}(\xi_{j_{1}}, \xi_{j_{2}}; e_{b}),
\end{equation}
exactly the right-hand side of clause~(ii) of the theorem.

Convergence of~\eqref{eq:mc5-g0g1-corner-trace-formula}: the sum over
the eval-gen-core basis is trace-class for $0 < |q| < 1$ by the
HS-sewing criterion applied to the genus-one component. The trace-class
property survives the residual genus-$0$ coupling because the residual
$2$-point function $\Phi^{\evcore}_{2}(\xi_{j_{1}}, \xi_{j_{2}}; e_{b})$
is bounded uniformly on compact subsets of the genus-$0$ residual sphere,
controlled by the conformal weight grading on $\cA^{\evcore}$.

\medskip
\textbf{Step 4: Compatibility with the bar differential at non-separating loci.}
The remaining content is the squaring-to-zero
$d_{\mathrm{bar}, 5}^{\, 2} = 0$ on the open stratum $\Conford_{5}(X)$
and its propagation across the boundary. By the climax theorem
\textup{(}Chapter~\ref{chap:climax-platonic}, equation~\eqref{eq:climax-equation-G2}\textup{)},
the bar differential is the pullback
\begin{equation}
\label{eq:mc5-g0g1-bar-differential-pullback}
d_{\mathrm{bar}, 5} \;=\; \KZ^{*}\bigl(\nabla^{(5)}_{\Arn}\bigr),
\end{equation}
and Arnold's three-term relation
\begin{equation}
\label{eq:mc5-g0g1-arnold-three-term}
\eta_{ij} \wedge \eta_{jk} + \eta_{jk} \wedge \eta_{ki} + \eta_{ki} \wedge \eta_{ij} = 0
\quad\text{on } \Conford_{5}(X), \quad 1 \le i < j < k \le 5,
\end{equation}
gives flatness $\nabla^{(5)\, 2}_{\Arn} = 0$. Pulling back, $d_{\mathrm{bar}, 5}^{\, 2} = 0$.
This is clause~(iii).

The Maurer-Cartan element $\tau^{\mathrm{sew}, 5}$ is then the universal
twisting morphism of the convolution complex
$\Conv(\BarOrd_{\le 5}(\cA^{\evcore}), \cA^{\evcore})$ at degree $5$, and
clause~(iv) asserts that it extends the MC4 element by construction:
each MC4 sewing prediction at $|S| = 2$ or $|S| = 3$ is a finite-stage
truncation of the MC5 element at $\le 4$ points, and the
genus-one clutching prediction on $\widehat{\Delta}^{(1)}_{T}$ is the
extension specified in clause~(ii).

\medskip
\textbf{Bar-cobar duality on the eval-gen-core.}
The completed counit $\widehat{\epsilon}_{5}$ is a quasi-isomorphism by
the strong-completion-tower closure of MC4
\textup{(}Theorem~\ref{thm:completed-bar-cobar-strong}\textup{)} applied
at finite stage $N = 5$, restricted to $\cA^{\evcore}_{\le 5}$. The
strong completion-tower hypothesis ensures the degree-cutoff lemma
\textup{(}Lemma~\ref{lem:degree-cutoff-bar-truncation}\textup{)} bounds
the bar differential by a finite sum at each weight stage, and
Mittag-Leffler is automatic on the eval-gen-core
\textup{(}Proposition~\ref{prop:standard-strong-filtration},
clause~(iv)\textup{)}. This proves clause~(v).

\medskip
\textbf{Cross-references.}
Clause~\textup{(i)} is the Beilinson--Drinfeld factorisation rule at
rational separating divisors~\cite{BeilinsonDrinfeld04}. Clause~\textup{(ii)}
is the genus-one trace factorisation on the clutching divisor
$\Delta^{(1)}_T$, with trace-class convergence supplied by the
HS-sewing criterion in the stated eval-gen-core ambient. The two
factorisations agree at the nodal limit $q=0$ because both reduce to the
same residue pairing through the node state.
\end{proof}

\section{Three clutching consequences}
\label{sec:mc5-g0g1-three-consequences}

The 5-point sewing theorem has three concrete instantiations, one for
each canonical landscape witness for which the elliptic supertrace in
clause~\textup{(ii)} is computable in closed form. We record them as
corollaries of
Theorem~\ref{thm:mc5-g0g1-wall-five-point-sewing}, with explicit scope
tags.

\begin{corollary}[Heisenberg elliptic function at the wall;
\ClaimStatusProvedHere]
\label{cor:mc5-g0g1-heisenberg-elliptic-function}
\index{Heisenberg!elliptic function at MC5 wall|textbf}
\textup{(UNIFORM-WEIGHT.)}
Let $\cA = H_{k}$ be the Heisenberg vertex algebra at level $k$, regarded
as a class-$\mathsf{G}$ chiral algebra in the standard landscape
\textup{(}Chapter~\ref{chap:shadow-quadrichotomy-platonic}\textup{)}. The
elliptic supertrace in clause~\textup{(ii)} of
Theorem~\ref{thm:mc5-g0g1-wall-five-point-sewing}, applied to
$\cA^{\evcore} = H_{k}$, is a Fredholm determinant
\begin{equation}
\label{eq:mc5-g0g1-heisenberg-fredholm-determinant}
F^{(1)}_{3}(\xi_{i_{1}}, \xi_{i_{2}}, \xi_{i_{3}}; e_{a}; q)
\;=\;
\det\bigl(1 - T_{q}\bigr)^{k} \cdot W_{3}(\xi_{i_{1}}, \xi_{i_{2}}, \xi_{i_{3}}; e_{a}; q),
\end{equation}
where $T_{q}$ is the genus-$1$ sewing kernel on the reduced Bergman
space and $W_{3}$ is the Wick-contraction polynomial of the three
insertions on the elliptic bubble. Combined with the residual $2$-point
function on the genus-$0$ component, the MC5 sewing reproduces the
genus-$1$ traced $5$-point function of the Heisenberg algebra.
\end{corollary}

\begin{proof}
For the Heisenberg algebra, Wick's theorem reduces the elliptic supertrace
to a sum of products of two-point functions on the elliptic bubble. The
infinite product structure comes from the bosonic Fock space character,
which contributes $\det(1 - T_{q})^{k}$ as the regularised product over
oscillator modes. The polynomial factor $W_{3}$ collects the connected
Wick contractions of the three insertions. This is exactly the
genus-one clutching chart of the preceding sewing theorem.
\end{proof}

\begin{corollary}[Virasoro class-$\mathsf{M}$ shadow corrections at the wall;
\ClaimStatusConditional]
\label{cor:mc5-g0g1-virasoro-class-m-corrections}
\index{Virasoro!shadow corrections at MC5 wall|textbf}
Let $\cA = \mathrm{Vir}_{c}$ be the Virasoro chiral algebra at generic
central charge $c$, regarded as a class-$\mathsf{M}$ chiral algebra in
the standard landscape \textup{(}Chapter~\ref{chap:shadow-quadrichotomy-platonic}\textup{)}.
The elliptic supertrace in clause~\textup{(ii)} of
Theorem~\ref{thm:mc5-g0g1-wall-five-point-sewing} acquires shadow
corrections at order $S_{r}$ for $r \ge 2$:
\begin{equation}
\label{eq:mc5-g0g1-virasoro-shadow-corrections}
F^{(1)}_{3}(\xi_{i_{1}}, \xi_{i_{2}}, \xi_{i_{3}}; e_{a}; q)
\;=\;
F^{(1), \mathrm{cl}}_{3}(\xi_{i_{1}}, \xi_{i_{2}}, \xi_{i_{3}}; e_{a}; q)
\;+\;
\sum_{r \ge 2} S_{r}(\mathrm{Vir}_{c}) \cdot
\delta^{(r)}_{3}(\xi_{i_{1}}, \xi_{i_{2}}, \xi_{i_{3}}; e_{a}; q)
\;+\;\delta F_{2}^{\mathrm{cross}},
\end{equation}
where $S_{r}(\mathrm{Vir}_{c})$ is the $r$-th Virasoro shadow coefficient
and $\delta^{(r)}_{3}$ is the corresponding $r$-loop shadow-Feynman
correction at the bubble. The term $\delta F_{2}^{\mathrm{cross}}$
records all-weight cross-channel mixing; it vanishes on the
uniform-weight Virasoro stress-tensor line. The genus-one clutching
chart is the first MC5 locus where these class-$\mathsf{M}$ corrections
are visible in the moduli geometry; the chain-level ambient question for the
resulting infinite series is treated in
Chapter~\ref{chap:mc5-class-m-chain-level-platonic}.
\end{corollary}

\begin{proof}
The Virasoro elliptic supertrace at $n = 3$ is computable as a series in
the bubble nome $q$ via the Virasoro character formula
$\chi_{\mathrm{Vir}_{c}}(q) = q^{-c/24} / \prod_{n \ge 2} (1 - q^{n})$.
The shadow corrections $S_{r}(\mathrm{Vir}_{c})$ enter as coefficients
of the curvature corrections $m_{0}^{\lfloor r/2 \rfloor - 1}$
in the bar complex
\textup{(}equation~\eqref{eq:harmonic-discrepancy} of
Chapter~\ref{chap:mc5-class-m-chain-level-platonic}\textup{)}, and the
sum~\eqref{eq:mc5-g0g1-virasoro-shadow-corrections} is the corresponding
shadow-Feynman expansion on the clutching chart. The chain-level
convergence of the resulting infinite series in the pro-ambient is the
content of
Theorem~\ref{thm:mc5-class-m-chain-level-pro-ambient}.
\end{proof}

\begin{corollary}[K3 elliptic genus at the wall; \ClaimStatusProvedHere]
\label{cor:mc5-g0g1-k3-elliptic-genus}
\index{K3 elliptic genus at MC5 wall|textbf}
Let $\cA = \cA(\mathrm{K3})$ be the chiral algebra of the K3 sigma model
at large volume, restricted to its sub-Sugawara stratum
$\{k < -h^{\vee}/2\}$ where the elliptic genus formula of
Eguchi--Ooguri--Tachikawa~\cite{EguchiOoguriTachikawa11} is unambiguously
defined as a Jacobi form of weight $0$ and index $1$. On the
genus-one component of the clutching chart
$\widehat{\Delta}^{(1)}_{T}$, the elliptic supertrace in
clause~\textup{(ii)} of
Theorem~\ref{thm:mc5-g0g1-wall-five-point-sewing} acquires the K3
elliptic genus as the bubble factor:
\begin{equation}
\label{eq:mc5-g0g1-k3-elliptic-genus-formula}
F^{(1)}_{3}(\xi_{i_{1}}, \xi_{i_{2}}, \xi_{i_{3}}; e_{a}; q, y)
\;=\;
\chi_{\mathrm{K3}}(q, y) \cdot W^{(1)}_{3, \mathrm{K3}}(\xi_{i_{1}}, \xi_{i_{2}}, \xi_{i_{3}}; e_{a}; q),
\end{equation}
up to a normalisation depending on the choice of insertion frames, where
$\chi_{\mathrm{K3}}(q, y)$ is the K3 elliptic genus and
$W^{(1)}_{3, \mathrm{K3}}$ is the connected $3$-point Wick polynomial on
the bubble in the K3 chiral algebra at the chosen stratum. MC5 thus
produces a chiral-algebraic derivation of the K3 elliptic genus from
the sewing principle, valid on the indicated sub-Sugawara stratum.
\end{corollary}

\begin{proof}
On the sub-Sugawara stratum $\{k < -h^{\vee}/2\}$, the K3 chiral algebra
admits an $\cN = 4$ superconformal structure with $\widehat{c} = 2$, and
the elliptic genus
\[
\chi_{\mathrm{K3}}(q, y)
\;=\;
\mathrm{Tr}_{\mathrm{RR}}\bigl((-1)^{F} y^{J_{0}} q^{L_{0} - c/24} \bar q^{\bar L_{0} - \bar c/24}\bigr)
\]
is a Jacobi form of weight $0$ and index $1$ in the variables $(q, y)$
\textup{(}see~\cite{EguchiOoguriTachikawa11}\textup{)}. The elliptic
supertrace on the clutching chart specialises to this elliptic genus
when the three insertions $(\xi_{i_{1}}, \xi_{i_{2}}, \xi_{i_{3}})$ are
chosen in the trivial conformal block at zero left-moving conformal
weight; in that case
$W^{(1)}_{3, \mathrm{K3}}(\xi_{i_{1}}, \xi_{i_{2}}, \xi_{i_{3}}; e_{a}; q)$
is a normalisation polynomial controlled by the Wick contractions of
the insertions on the elliptic bubble. The detailed identification of
the bubble factor with $\chi_{\mathrm{K3}}(q, y)$ uses the
Eguchi--Ooguri--Tachikawa decomposition of $\chi_{\mathrm{K3}}$ into
$\cN = 4$ characters, applied to the trace
formula~\eqref{eq:mc5-g0g1-genus-1-traced-three-point}. The scope is
restricted to the sub-Sugawara stratum because the elliptic-genus
factorisation through this specific chiral algebra is unambiguous
there. The elliptic genus itself is deformation invariant; what is not
claimed on the full Sugawara stratum is this MC5 chiral-algebraic
factorisation of the genus-$1$ bubble trace.
\end{proof}

\begin{remark}[Sub-Sugawara comparison locus]
\label{rem:mc5-g0g1-k3-sub-sugawara-scope}
Corollary~\ref{cor:mc5-g0g1-k3-elliptic-genus} is deliberately
restricted to the sub-Sugawara comparison locus. The K3 elliptic genus
is a fixed weak Jacobi form of weight~$0$ and index~$1$; Mathieu
moonshine concerns its character decomposition, not a moduli-dependent
correction to the genus itself.
The sub-Sugawara stratum $\{k < -h^{\vee}/2\}$ is the regime where the
chiral 5-point sewing of MC5 reduces to the elementary Jacobi-form
formula of Eguchi--Ooguri--Tachikawa. Outside this stratum, the same
elliptic genus may still be present, but MC5 does not by itself supply
the required chiral trace model. The full-moduli extension is not
claimed here.
\end{remark}

\section{The closed MC1 -- MC5 chain}
\label{sec:mc5-g0g1-closed-mc-chain}

With Theorem~\ref{thm:mc5-g0g1-wall-five-point-sewing} in place, the
MC chain reads as a sequence of sewing theorems for the $n$-point
compactifications. MC5 adds the first genus-one clutching datum to the
rational genus-zero boundary calculus.

\begin{itemize}
\item \textbf{MC1.} Augmentation: the MC element at $n = 1$ is the
unit-vacuum identification.
\item \textbf{MC2.} Binary OPE: at $n = 2$, the MC element is the
collision residue $r(z) = \mathrm{Res}^{\mathrm{coll}}_{0, 2}(\Theta_{\cA})$,
the chiral classical $r$-matrix.
\item \textbf{MC3.} Associator: at $n = 3$, the MC element on the
eval-gen-core encodes the Stasheff pentagon at the level of the bar
coderivation $b_{2} + b_{3}$
\textup{(}Theorem~\ref{thm:mc3-evaluation-core-five-family}\textup{)}.
\item \textbf{MC4.} Quadrilateral sewing: at $n = 4$, the MC element
extends across the boundary $\partial \Mbar_{0,4}$ via the
strong-completion-tower closure on the eval-gen-core
\textup{(}Theorem~\ref{thm:completed-bar-cobar-strong}\textup{)}.
\item \textbf{MC5.} Genus-one clutching: at $n = 5$, the MC element
extends across the rational separating divisors $\delta_{0|S, S^c}$
\textup{(}governed by MC4\textup{)} and is compatible with the clutching
charts $\widehat{\Delta}^{(1)}_{T}$ \textup{(}governed by the elliptic
$3$-point trace on the genus-one component\textup{)}, restricted to the eval-gen-core
\textup{(}Theorem~\ref{thm:mc5-g0g1-wall-five-point-sewing}\textup{)}.
\end{itemize}

\begin{remark}[Climax-theorem integration]
\label{rem:mc5-g0g1-climax-integration}
Theorem~\ref{thm:mc5-g0g1-wall-five-point-sewing} is the $n = 5$
specialisation of universal KZ sewing
\textup{(}Chapter~\ref{chap:climax-platonic},
equation~\eqref{eq:climax-equation-G2}\textup{)}. Pulling back
$\nabla^{(5)}_{\Arn}$ from $\Conford_{5}(X)$ to its compactification
$\widehat{\Mbar}_{0,5}(X)$ produces a flat connection on the boundary
divisors; the pullback agrees with the chiral bar differential
$d_{\mathrm{bar}, 5}$ by clause~(iii), and Arnold's three-term relation
\eqref{eq:mc5-g0g1-arnold-three-term} becomes the squaring-to-zero
$d_{\mathrm{bar}, 5}^{\, 2} = 0$.

On the genus-one clutching chart one replaces the rational Arnold forms
by the elliptic KZB forms, equivalently by Felder's elliptic
$r$-matrix on the genus-one component. Clause~\textup{(ii)} is the
corresponding elliptic trace block, glued to the residual rational
two-point block through the node pairing.

The per-family $\kappa$ conductor of the harmonic discrepancy
\textup{(}Chapter~\ref{chap:mc5-class-m-chain-level-platonic},
equation~\eqref{eq:harmonic-discrepancy}\textup{)} is the negative of
the ghost charge of any quasi-free BRST resolution by the universal
identity of Chapter~\ref{chap:universal-conductor}. On the genus-one
clutching chart, this $\kappa$ conductor governs the magnitude of the
shadow corrections of
Corollary~\ref{cor:mc5-g0g1-virasoro-class-m-corrections}.

By the universality~\eqref{eq:climax-equation-G2}, any other $5$-point
chiral sewing prescription \textup{(}Drinfeld--Kohno, Verlinde,
Borcherds, Felder\textup{)} is a pullback of $\nabla^{(5)}_{\Arn}$ along
the structure functor $\cA \to \mathrm{End}(\BarOrd(\cA))$ on the
rational genus-zero surface. The theorem applies to those standard
landscape objects for which the eval-gen-core and trace-class
hypotheses are verified.
\end{remark}

\section{The Clutching Limit}
\label{sec:mc5-g0g1-inner-motion}

Fix $T=\{i_1,i_2,i_3\}$. The clutching map
\eqref{eq:mc5-genus-one-clutching} glues a genus-one component carrying
the markings in~$T$ to a genus-zero component carrying the two remaining
markings. The smoothing parameter $q$ belongs to the genus-one
component; the nodal limit $q=0$ recovers the ordinary residue pairing
through the gluing state.

Theorem~\ref{thm:mc5-g0g1-wall-five-point-sewing} certifies that the
genus-zero rational extension and the genus-one clutching block are
compatible at that nodal limit. On the rational side, the correlator is
the holomorphic $5$-point function on $\Conford_{5}(X)$ extended across
the ten divisors of $\Mbar_{0,5}$. On the genus-one side, the trace block
is the elliptic $3$-point function on the genus-one component paired
with the residual genus-zero $2$-point function. Clause~\textup{(ii)}
is exactly this pairing.

Every higher MC stage must impose analogous compatibility along the
relevant clutching morphisms. The chain-level companion
\textup{(}Chapter~\ref{chap:mc5-class-m-chain-level-platonic}\textup{)}
addresses the ambient-completion question for class $\mathsf{M}$ chiral
algebras in the pro-ambient.

\section{Summary}
\label{sec:mc5-g0g1-summary}

The MC5 statement is the following.

\begin{quote}
\emph{Let $X$ be a smooth projective curve and let $\cA$ be a strong
completion tower of finite resonance rank, restricted to its
evaluation-generated core $\cA^{\evcore}$. The $5$-point chiral
correlator of $\cA^{\evcore}$ on $X$ extends uniquely over the completed
genus-zero compactification $\widehat{\Mbar}_{0,5}(X)$ and is compatible
with the genus-one clutching charts
$\widehat{\Delta}^{(1)}_{T}$ for $|T|=3$. The rational extension is
governed at the ten separating divisors by the MC4 sewing prediction.
The clutching block is the trace of the genus-one $3$-point function
against the residual genus-zero $2$-point function. The bar differential
is the pullback
of Arnold's universal flat connection along the KZ functor of
Chapter~\ref{chap:climax-platonic}; flatness gives squaring-to-zero of
the bar differential, and the resulting Maurer-Cartan element extends
the MC4 element. The completed bar-cobar counit at finite stage
$N = 5$ is a quasi-isomorphism. The eval-core qualifier restricts the
theorem to evaluation-generated representations; full-category
extension to $\DK_g$ is conjectural
\textup{(}Conjecture~\ref{conj:mc5-g0g1-wall-full-category}\textup{)}.}
\end{quote}

The MC chain $\mathrm{MC1} \to \mathrm{MC2} \to \mathrm{MC3} \to
\mathrm{MC4} \to \mathrm{MC5}$ thus reaches its first genus-one
clutching compatibility, with each MC stage being the sewing theorem at
its $n$-point level. Three concrete consequences (the Heisenberg
elliptic Fredholm determinant, the Virasoro class-$\mathsf{M}$ shadow
corrections, the K3 elliptic genus on the sub-Sugawara stratum)
illustrate the theorem on the canonical landscape witnesses, with the
uniform-weight scope made explicit and the sub-Sugawara restriction in
the K3 case.

\begin{center}
\rule{0.3\textwidth}{0.4pt}
\end{center}

The MC5 surface has two components. The geometric $5$-point sewing with
genus-one clutching is proved on the eval-gen-core
\textup{(}Theorem~\ref{thm:mc5-g0g1-wall-five-point-sewing}\textup{)};
the chain-level chain-ambient companion for class $\mathsf{M}$ is proved
in the pro-ambient \textup{(}Chapter~\ref{chap:mc5-class-m-chain-level-platonic}\textup{)};
together they constitute the closed MC5 statement. Higher-stage
extensions are conjectural
\textup{(}Remark~\ref{rem:mc5-g0g1-higher-stage-conjectural}\textup{)}.
